\section{Related Work} \label{sec:related-work}
  % Rowhammer in General
  It all began in 2014, when Kim et al.~\cite{paperFlippingBits} discovered the rowhammer vulnerability 
  in the former DDR3 of an x86 system and talks about their technical details.
  Seaborn et al.~\cite{paperExploitingTheDRAM} subsequently developed a working  probabilistic  exploit 
  that use this effect for local privilege escalation on year later.
  Shortly thereafter, in 2016, Gruss et al.~\cite{paperRowhammerjs} demonstrated a Rowhammer attack 
  in JavaScript that did not rely on the \texttt{CLFLUSH} instruction, instead 
  evicting cache lines from the CPU cache~\cite{paperRowhammerOnSteroids}. This work showed that the previously 
  considered essential instruction was no longer required and revealed a wide 
  range of additional exploitation possibilities of the vulnerability.
  In the same year, Pessl et al.~\cite{paperDRAMA} introduced an automated reverse-engineering 
  approach to identify DRAM addressing functions using their tool~\cite{toolDrama}. Their 
  method leveraged timing-based measurements to systematically reconstruct the 
  mapping of physical addresses to DRAM channels, ranks, and banks, thereby enabling 
  more precise and targeted memory-based attacks.
  Schwarz et al.~\cite{thesisDRAMA} extended the theoretical and practical foundations of 
  DRAMA introduced by Pessl et al., adapting the approach for ARM-based systems. 
  Although the original implementation was not publicly available, their work 
  provided the basis for our thesis to port DRAMA to ARM, enabling systematic 
  analysis of DRAM buffer on modern ARM platforms.
  \\
  % Rowhammer on ARM platforms
  Earlier Rowhammer attacks primarily targeted x86 systems, largely because 
  the unprivileged \texttt{CLFLUSH} instruction enables efficient 
  exploitation~\cite{paperFlippingBits, paperRHonRi3B+}. As ARM has become the dominant architecture 
  in mobile devices, researchers have attempted to adapt these x86-oriented 
  techniques. However, this translation is challenging, and some have even 
  suggested that Rowhammer attacks on ARM might be infeasible, partly due 
  to the potentially slower memory controllers that could prevent the bug 
  from being triggered~\cite{paperDrammer,paperTriggeringRHHardwareFault}.
  Nevertheless, subsequent studies demonstrated that Rowhammer can indeed be 
  exploited on ARM by identifying alternative methods to bypass these 
  architectural limitations. \cite{paperDrammer} were the first study of Rowhammer prevalence 
  on ARM-based devices, including Android smartphones and Chromebooks, as it 
  was released from Van der Veen et al. in 2016. They provide a deterministic 
  Rowhammer attack that cannot be mitigated by current defenses by leveraging 
  user-accessible DMA buffers on Android devices, which allow memory accesses 
  to bypass the CPU caches and directly interact with DRAM.
  Lipp et al.~\cite{paperARMageddon} demonstrated in 2016 cache side-channel attacks on ARM 
  based mobile devices without requiring privileged instructions and to 
  replace \texttt{CLFLUSH}, enabling fine-grained timing attacks on ARM.
  Subsequently, the authors of~\cite{thesisCacheAttacksandRH} extended such cache attacks on ARM 
  architectures to practical exploitation scenarios, including keystroke 
  inference, attacks on cryptographic algorithms, and the induction of 
  Rowhammer faults. This methodology serves as a foundation in our thesis, 
  where it is adopted as one approach to effectively bypass the CPU cache.
  Zhang et al.~\cite{paperTriggeringRHHardwareFault} validates in 2018 the attack on ARMv8-based single-board 
  computers, namely by Amlogic S905X SBC by reliably trigger bitflips through 
  ARM cache maintenance instructions analogous to the case of x86~\cite{paperRHonRi3B+}.
  \\
  % Previous attacks on Raspberry Pi's
  Bekele et al.~\cite{paperRHonRi3B+} demonstrated the susceptibility of the Raspberry Pi 3B+ 
  to Rowhammer attacks by leveraging ARM cache maintenance instructions. In 
  addition, they analyzed the impact of DSB on attack effectiveness and examined 
  the likelihood of inducing repeated bitflips at specific memory locations.
  Plin et al.~\cite{paperKnockKnock} introduced in 2025 a black-box approach to reverse-engineer 
  DRAM address mappings without relying on proprietary hardware documentation, 
  evaluating their method on a comparable platform, namely the Raspberry Pi 4. 
  Nevertheless, they observed and induced a closed-page policy on this system, influencing 
  the effectiveness of reconstructing the underlying DRAM address mapping dramatically.
  \\
  % Research Gap
  Previous studies have demonstrated the susceptibility of Rowhammer on older 
  Raspberry Pi models. Even so, no other embedded systems within the Raspberry Pi 
  series were evaluated. To date, research on Rowhammer has been limited almost 
  exclusively to the Raspberry Pi 4, while the Raspberry Pi 5 remains entirely 
  unexplored. Whereas the Raspberry Pi 4 has only been scarcely investigated and 
  explored to trigger Rowhammer faults with \gls{gpu}-based memory access patterns~\cite{paperOpenGL} 
  in 2025.
  Consequently, it remains an open question whether modern Raspberry Pi platforms, 
  such as the Raspberry Pi 4 and Raspberry Pi 5, are susceptible to software-induced 
  Rowhammer attacks.
